% PDF transcription; source and license in ../../tex/005_半单Lie代数表示的完全可约性.tex
\paragraph{18.1. Extensions.}
Let $\mathfrak g$ be a Lie algebra and $U,W$ be representations of $\mathfrak g$. We would like to classify all representations $V$ which fit into a short exact sequence
\[
0\longrightarrow U\longrightarrow V\longrightarrow W\longrightarrow0,\tag{18.1}
\]
i.e., $U\subset V$ is a subrepresentation such that the surjection $p:V\to W$ has kernel $U$ and thus defines an isomorphism $V/U\cong W$. In other words, $V$ is endowed with a 2-step filtration with $F_0V=U$ and $F_1V=V$ such that $F_1V/F_0V=W$, so $\operatorname{gr}(V)=U\oplus W$. To do so, pick a splitting of this sequence as a sequence of vector spaces, i.e. an injection $i:W\to V$ (not a homomorphism of representations, in general) such that $p\circ i=\operatorname{Id}_W$. This defines a linear isomorphism $\widetilde i:U\oplus W\to V$ given by $(u,w)\mapsto u+i(w)$, which allows us to rewrite the action of $\mathfrak g$ on $V$ as an action on $U\oplus W$. Since $\widetilde i$ is not in general a morphism of representations, this action is given by
\[
\rho(x)(u,w)=(xu+a(x)w,xw)
\]
where $a:\mathfrak g\to\Hom_k(W,U)$ is a linear map, and $\widetilde i$ is a morphism of representations iff $a=0$.

What are the conditions on $a$ to give rise to a representation? We compute:
\[
\begin{aligned}
\rho([x,y])(u,w)&=([x,y]u+a([x,y])w,[x,y]w),\\
[\rho(x),\rho(y)](u,w)&=([x,y]u+([x,a(y)]+[a(x),y])w,[x,y]w).
\end{aligned}
\]
Thus the condition to give a representation is the Leibniz rule
\[
a([x,y])=[x,a(y)]+[a(x),y]=[x,a(y)]-[y,a(x)].
\]
In general, if $E$ is a representation of $\mathfrak g$ then a linear function $a:\mathfrak g\to E$ such that
\[
a([x,y])=x\circ a(y)-y\circ a(x)
\]
is called a \emph{1-cocycle} of $\mathfrak g$ with values in $E$. The space of 1-cocycles is denoted by $Z^1(\mathfrak g,E)$.

\paragraph{Example 18.1.}
We have $Z^1(\mathfrak g,k)=(\mathfrak g/[\mathfrak g,\mathfrak g])^*$ and $Z^1(\mathfrak g,\mathfrak g)=\operatorname{Der}\mathfrak g$.

Thus we see that in our setting $a:\mathfrak g\to\Hom_k(W,U)$ defines a representation if and only if $a\in Z^1(\mathfrak g,\Hom_k(W,U))$. Denote the representation $V$ attached to such $a$ by $V_a$. Then we have a natural short exact sequence
\[
0\to U\to V_a\to W\to0.
\]
It may, however, happen that some $a\ne0$ defines a trivial extension $V\cong U\oplus W$, i.e., $V_a\cong V_0$, and more generally $V_a\cong V_b$ for $a\ne b$.

Let us determine when this happens. More precisely, let us look for isomorphisms $f:V_a\to V_b$ preserving the structure of the short exact sequences, i.e., such that $\operatorname{gr}(f)=\operatorname{Id}$. Then
\[
f(u,w)=(u+Aw,w)
\]
where $A:W\to U$ is a linear map. Then we have
\[
xf(u,w)=x(u+Aw,w)=(xu+xAw+b(x)w,xw)
\]
and
\[
fx(u,w)=f(xu+a(x)w,xw)=(xu+a(x)w+Axw,xw),
\]
so we get that $xf=fx$ iff
\[
[x,A]=a(x)-b(x).
\]
In particular, setting $b=0$, we see that $V$ is a trivial extension if and only if $a(x)=[x,A]$ for some $A$.

More generally, if $E$ is a $\mathfrak g$-module, the linear function $a:\mathfrak g\to E$ given by $a(x)=xv$ for some $v\in E$ is called the \emph{1-coboundary} of $v$, and one writes $a=dv$. The space of 1-coboundaries is denoted by $B^1(\mathfrak g,E)$; it is easy to see that it is a subspace of $Z^1(\mathfrak g,E)$, i.e., a 1-coboundary is always a 1-cocycle. Thus in our setting $f:V_a\to V_b$ is an isomorphism of representations iff
\[
a-b=dA,
\]
i.e., there is an isomorphism $f:V_a\cong V_b$ with $\operatorname{gr}(f)=\operatorname{Id}$ if and only if $a=b$ in the quotient space
\[
\operatorname{Ext}^1(W,U):=Z^1(\mathfrak g,\Hom_k(W,U))/B^1(\mathfrak g,\Hom_k(W,U)).
\]
The notation is justified by the fact that this space parametrizes extensions of $W$ by $U$. More precisely, every short exact sequence (18.1) gives rise to a class $[V]\in\operatorname{Ext}^1(W,U)$, and the extension defined by this sequence is trivial iff $[V]=0$.

More generally, for a $\mathfrak g$-module $E$ the space
\[
H^1(\mathfrak g,E):=Z^1(\mathfrak g,E)/B^1(\mathfrak g,E)
\]
is called the \emph{first cohomology} of $\mathfrak g$ with coefficients in $E$. Thus,
\[
\operatorname{Ext}^1(W,U)=H^1(\mathfrak g,\Hom_k(W,U)).
\]

\begin{lemma}[18.2; standard exactness statement used below]
A short exact sequence $0\to U\to V\to W\to0$ gives rise to an exact sequence
\[
H^1(\mathfrak g,U)\longrightarrow H^1(\mathfrak g,V)\longrightarrow H^1(\mathfrak g,W).
\]
\end{lemma}
\noindent\emph{Exercise 18.3.} Prove Lemma 18.2.

\paragraph{18.2. Whitehead's theorem.}
We have shown in Corollary 17.6 and Proposition 17.9 that for a semisimple $\mathfrak g$ over a field of characteristic zero, $H^1(\mathfrak g,k)=(\mathfrak g/[\mathfrak g,\mathfrak g])^*=0$, and $H^1(\mathfrak g,\mathfrak g)=\operatorname{Der}\mathfrak g/\mathfrak g=0$. In fact, these are special cases of a more general theorem.

\begin{theorem}[18.4; Whitehead]
If $\mathfrak g$ is semisimple in characteristic zero then for every finite dimensional representation $V$ of $\mathfrak g$, $H^1(\mathfrak g,V)=0$.
\end{theorem}

\paragraph{18.3. Proof of Theorem 18.4.}
We will use the following lemma, which holds over any field.
\begin{lemma}[18.5]
Let $E$ be a representation of a Lie algebra $\mathfrak g$ and $C\in U(\mathfrak g)$ be a central element which acts by $0$ on the trivial representation of $\mathfrak g$ and by some scalar $\lambda\ne0$ on $E$. Then $H^1(\mathfrak g,E)=0$.
\end{lemma}
\begin{proof}
We have seen that $H^1(\mathfrak g,E)=\operatorname{Ext}^1(k,E)$, so our job is to show that any extension
\[
0\to E\to V\to k\to0
\]
splits. Let $p:V\to k$ be the projection. We claim that there exists a unique vector $v\in V$ such that $p(v)=1$ and $Cv=0$. Indeed, pick some $w\in V$ with $p(w)=1$. Then $Cw\in E$, so set $v=w-\lambda^{-1}Cw$. Since $C^2w=\lambda Cw$, we have $Cv=0$. Also if $v'$ is another such vector then $v-v'\in E$ so $C(v-v')=\lambda(v-v')=0$, hence $v=v'$.

Thus $kv\subset V$ is a $\mathfrak g$-invariant complement to $E$ (as $C$ is central), which implies the statement.
\end{proof}

It remains to construct a central element of $U(\mathfrak g)$ for a semisimple Lie algebra $\mathfrak g$ to which we can apply Lemma 18.5. This can be done as follows. Let $a_i$ be a basis of $\mathfrak g$ and $a^i$ the dual basis under an invariant inner product on $\mathfrak g$ (for example, the Killing form). Define the \emph{(quadratic) Casimir element}
\[
C:=\sum_i a_i a^i.
\]
It is easy to show that $C$ is independent on the choice of the basis (although it depends on the choice of the inner product). Also $C$ is central: for $y\in\mathfrak g$,
\[
[y,C]=\sum_i([y,a_i]a^i+a_i[y,a^i])=0
\]
since
\[
\sum_i([y,a_i]\otimes a^i+a_i\otimes[y,a^i])=0
\]
(this is seen by taking the inner product of the first tensorand with $a^j$ and using the invariance of the inner product). Finally, note that for $\mathfrak g=\mathfrak{sl}_2$, $C$ is proportional to the Casimir element $2fe+\frac{h^2}{2}+h=ef+fe+\frac{h^2}{2}$ considered previously, as the basis $f,e,\frac h{\sqrt2}$ is dual to the basis $e,f,\frac h{\sqrt2}$ under an invariant inner product of $\mathfrak g$.

The key lemma used in the proof of Theorem 18.4 is the following.
\begin{lemma}[18.6]
Let $\mathfrak g$ be semisimple in characteristic zero and $V$ be a nontrivial finite dimensional irreducible $\mathfrak g$-module. Then there is a central element $C\in U(\mathfrak g)$ such that $C|_k=0$ and $C|_V\ne0$.
\end{lemma}
\begin{proof}
Consider the invariant symmetric bilinear form on $\mathfrak g$
\[
B_V(x,y)=\operatorname{Tr}|_V(xy).
\]
We claim that $B_V\ne0$. Indeed, let $\overline{\mathfrak g}\subset\mathfrak{gl}(V)$ be the image of $\mathfrak g$. By Lemma 17.1, if $B_V=0$ then $\overline{\mathfrak g}$ is solvable, so, being the quotient of a semisimple Lie algebra $\mathfrak g$, it must be zero, hence $V$ is trivial, a contradiction.

Let $I=\ker(B_V)$. Then $I\subset\mathfrak g$ is an ideal, so by Proposition 17.7, $\mathfrak g=I\oplus\mathfrak g'$ for some semisimple Lie algebra $\mathfrak g'$, and $B_V$ is nondegenerate on $\mathfrak g'$. Let $C$ be the Casimir element of $U(\mathfrak g')$ corresponding to the inner product $B_V$. Then $\operatorname{Tr}_V(C)=\sum_i B_V(a_i,a^i)=\dim\mathfrak g'$, so $C|_V=\frac{\dim\mathfrak g'}{\dim V}\ne0$. Also it is clear that $C|_k=0$, so the lemma follows.
\end{proof}

\begin{corollary}[18.7]
For any irreducible finite dimensional representation $V$ of a semisimple Lie algebra $\mathfrak g$ over a field $k$ of characteristic zero, we have $H^1(\mathfrak g,V)=0$.
\end{corollary}
\begin{proof}
If $V$ is nontrivial, this follows from Lemmas 18.5 and 18.6. On the other hand, if $V=k$ then $H^1(\mathfrak g,V)=(\mathfrak g/[\mathfrak g,\mathfrak g])^*=0$.
\end{proof}

Now we can prove Theorem 18.4. By Lemma 18.2, it suffices to prove the theorem for irreducible $V$, which is guaranteed by Corollary 18.7.

\paragraph{18.4. Complete reducibility of representations of semisimple Lie algebras.}
\begin{theorem}[18.9]
Every finite dimensional representation of a semisimple Lie algebra $\mathfrak g$ over a field of characteristic zero is completely reducible, i.e., isomorphic to a direct sum of irreducible representations.
\end{theorem}
\begin{proof}
Theorem 18.4 implies that for any finite dimensional representations $W,U$ of $\mathfrak g$ one has $\operatorname{Ext}^1(W,U)=0$. Thus any short exact sequence
\[
0\to U\to V\to W\to0
\]
splits, which implies the statement.
\end{proof}
